\documentclass[11pt]{article}

\usepackage[margin=1in]{geometry}
\usepackage{amsmath,amssymb}
\usepackage[T1]{fontenc}
\usepackage[utf8]{inputenc}
\usepackage{lmodern}

\title{Rewritten Reasoning Log for the Closed-Form Two-Minus Water-Wave Amplitude}
\author{}
\date{}

\begin{document}
\maketitle


\section*{Rewritten Reasoning Log}

The problem is to identify a closed-form analytic expression for the tree-level
on-shell water-wave amplitude in the sector
\[
    \sigma=(-1,-1,+1,\ldots,+1),
\]
with all frequencies incoming and with momenta fixed by
\[
    k_i=\sigma_i\omega_i^2/g.
\]
On the resonant manifold the two conservation laws are
\[
    \sum_{i=1}^n \omega_i=0,
    \qquad
    \sum_{i=1}^n \sigma_i\omega_i^2=0.
\]
The supplied file \texttt{OnShellBG.m} contains an exact Berends--Giele
implementation. The practical goal was therefore not to rederive the whole
recursion from first principles, but to use exact generated data from that
implementation to find the chamber structure and the polynomial that appears
inside each chamber.

The first step was to isolate the reusable definitions from
\texttt{OnShellBG.m}. A generated helper file, \texttt{bg\_core.wl}, copied the
definitions for the kernels, vertex, propagator, set partitions, BG current,
amplitude, and kinematic solver, without running the unrelated sample tests at
the bottom of the original file. This made repeated rational evaluations
straightforward. All subsequent exploratory files were generated inside the
case directory and then moved into \texttt{case\_2/codex\_55\_xhigh}.

Exact sample evaluations showed the basic shape of the answer. For instance,
with five points and free frequencies \(\{2,5/2,3\}\), the solver gives
\[
    \omega=\{-9/2,\,2,\,5/2,\,3,\,-3\}
\]
and the BG code gives
\[
    A_5=-2304\,i.
\]
Other five-point chambers gave different homogeneous polynomials, for example
\[
\begin{array}{c|c|c}
\text{free frequencies} & \omega & A_5/i\\
\hline
\{5,1,2\} & \{-11/4,5,1,2,-21/4\} & -1760\\
\{-1,2,5\} & \{-16/3,-1,2,5,-2/3\} & 14336/243\\
\{1,-2,5\} & \{-11/2,1,-2,5,3/2\} & -88
\end{array}
\]
The amplitude was consistently purely imaginary in these checks.

A scaling test fixed the expected homogeneous degree. Scaling the five-point
sample by \(\lambda\) gave ratios
\[
    1,\quad 64,\quad 729,\quad 1/64
\]
for \(\lambda=1,2,3,1/2\). Thus the five-point amplitude scales like
\(\lambda^6\), matching the pattern \(2n-4\). This was useful because it
constrained the possible polynomial form and made it clear that the observed
piecewise expressions were not accidental rational functions in the
frequencies.

The four-point case required separate care. Real two-minus four-point
resonance forces pairwise cancellations in the momenta, and the unregularized
BG recursion can encounter internal zero-momentum currents, producing
\(\texttt{Indeterminate}\). The finite value was checked by splitting the two
positive external momenta by a symbolic parameter \(\delta\), preserving total
momentum, evaluating the BG expression, and taking \(\delta\to0^+\). For
example,
\[
    \omega=\{-3,2,3,-2\}
\]
gave
\[
    \lim_{\delta\to0^+} A_4/i=-192.
\]
This matched the eventual closed form and explained why raw four-point calls
needed a limiting interpretation.

Several structural checks helped narrow the formula. A cubic-only BG recursion
was tested as a possible simplification, but it did not match the full answer:
for the five-point sample \(\{2,5/2,3\}\), the full value was
\(-2304\,i\), while the cubic-only value was \(-5296\,i\). Contact terms
therefore matter. Kernel probes also showed sharp sign behavior: with all
positive momenta, the generated \(E\) and \(F\) kernels agreed in small cases,
while representative mixed-sign inputs vanished. This suggested that the
absolute values in the kernels were the source of the chamber decomposition.

The decisive computation was sign-resolved symbolic evaluation. In a fixed
sample chamber, every \(\operatorname{Abs}\) in the BG expression can be
replaced by either \(+x\) or \(-x\), using the sample point to choose the sign.
For five points with symbolic free variables \(x,y,z\), the on-shell solution is
\[
    \omega_5=-\frac{(x+y)(x+z)}{x+y+z},
    \qquad
    \omega_1=-(x+y+z+\omega_5).
\]
In the chamber represented by \(\{x,y,z\}=\{2,5/2,3\}\), the exact symbolic
BG result reduced to
\[
    \frac{A_5}{i}
    =
    \frac{-16x^5(xy+y^2+xz+yz+z^2)}{x+y+z}.
\]
After substituting the on-shell expression for \(\omega_1\), this is simply the
corresponding chamber polynomial in the physical frequencies. Other chambers
produced different-looking expressions, but after rewriting them in terms of
the squared positive-leg frequencies and the smaller of the two negative
momentum magnitudes, the same inclusion--exclusion pattern appeared.

Here is the more explicit bridge from the five-point data to the
inclusion--exclusion ansatz. At five points I normalized away the universal
factor
\[
    i\,16\,\omega_1\omega_2
\]
and studied
\[
    G_5=\frac{A_5}{i\,16\,\omega_1\omega_2}.
\]
The first symbolic chamber, represented by
\[
    \{x,y,z\}=\{2,5/2,3\},
    \qquad
    \omega=\{-9/2,2,5/2,3,-3\},
\]
has
\[
    r=\min(\omega_1^2,\omega_2^2)=\omega_2^2=x^2.
\]
All three positive-leg squares are larger than \(r\):
\[
    \omega_3^2=25/4,\qquad \omega_4^2=9,\qquad \omega_5^2=9,
\]
so no nonempty subset of positive legs can be active. The symbolic result
\[
    \frac{A_5}{i}=16\omega_1\omega_2\,r^2
\]
therefore looks like the empty-subset term of a truncated expression.

The next useful chamber was
\[
    \omega=\{-11/4,5,1,2,-21/4\}.
\]
Here
\[
    r=\omega_1^2=\frac{121}{16},
    \qquad
    q_3=1,\quad q_4=4,\quad q_5=\frac{441}{16}.
\]
Only the subsets of \(\{3,4\}\) fit below \(r\); anything containing leg \(5\)
is inactive. Dividing the BG value by \(i\,16\omega_1\omega_2\) gives
\[
    G_5=8.
\]
The finite-difference combination
\[
    r^2-(r-q_3)^2-(r-q_4)^2+(r-q_3-q_4)^2
\]
also equals \(8\). In fact, when \(r>q_3+q_4\), this expression collapses to
\[
    2q_3q_4.
\]
This matches the sign-resolved symbolic formula in that chamber,
\[
    \frac{A_5}{i}=32\,\omega_1\omega_2\,q_3q_4.
\]

A chamber with only one nonempty active subset gave the same clue. For
\[
    \omega=\{-16/3,-1,2,5,-2/3\}
\]
one has
\[
    r=1,
    \qquad
    q_3=4,\quad q_4=25,\quad q_5=\frac{4}{9}.
\]
Only \(q_5\) is below \(r\), so the predicted normalized expression is
\[
    G_5=r^2-(r-q_5)^2
       =1-\left(1-\frac{4}{9}\right)^2
       =\frac{56}{81}.
\]
Multiplying back by \(16\omega_1\omega_2=256/3\) gives
\[
    \frac{A_5}{i}=\frac{256}{3}\,\frac{56}{81}
    =\frac{14336}{243},
\]
which is exactly the BG value.

These examples are the point where the pattern became more specific than
``some piecewise degree-six polynomial.''  The normalized five-point answer was
not depending on the signed frequencies separately, except through
\(\omega_1\omega_2\), the cutoff \(r\), and the positive-leg squares \(q_j\).
The active terms were controlled by subset-sum inequalities \(Q_S<r\), and
their signs were exactly the alternating signs of an inclusion--exclusion
finite difference of \(t^2\):
\[
    G_5
    =
    \sum_{S\subseteq\{3,4,5\}}
       (-1)^{|S|}(r-Q_S)_+^2.
\]
The notation \((\cdot)_+\) is what turns the chamber-dependent list of active
subsets into one compact expression.

Let
\[
    P=\{3,4,\ldots,n\},\qquad
    q_j=\omega_j^2\quad(j\in P),
\]
and set
\[
    r=\min(\omega_1^2,\omega_2^2),
    \qquad
    Q_S=\sum_{j\in S}q_j.
\]
Define the truncated power by
\[
    (x)_+^m=
    \begin{cases}
    x^m, & x>0,\\
    0, & x<0.
    \end{cases}
\]
On chamber boundaries the continuous limiting value is used. The resulting
formula is
\[
    A_n
    =
    i\,2^{n-1}\omega_1\omega_2
    \sum_{S\subseteq\{3,\ldots,n\}}
       (-1)^{|S|}\bigl(r-Q_S\bigr)_+^{\,n-3}.
\]
The chambers are exactly the regions where the signs of
\[
    \omega_1^2-\omega_2^2
    \quad\text{and}\quad
    r-Q_S\quad(S\subseteq\{3,\ldots,n\})
\]
are fixed. Inside any one chamber, the inactive terms with \(Q_S>r\) are
removed, and the remaining finite sum is an ordinary homogeneous polynomial of
degree \(2n-4\).

This expression also explains the low-point observations. At five points the
power is \(2\), and the active subset pattern selects the chamber polynomial.
At six and seven points, the same normalized expression appears with powers
\(3\) and \(4\). Adding one positive leg has the finite-difference form
\[
    F(r;q_1,\ldots,q_m)
    =
    F(r;q_1,\ldots,q_{m-1})
    -
    F(r-q_m;q_1,\ldots,q_{m-1}),
\]
whose closed solution is precisely the subset sum above. This is the compact
reason that the formula is piecewise polynomial rather than a single global
polynomial.

The final verification was exact rational comparison with the supplied BG code.
For non-boundary cases the script \texttt{verify.wl} evaluated both
\texttt{BGAmplitude} and the closed form. The reported difference was exactly
zero in all packaged checks:
\[
\begin{array}{c|c|c|c}
n & \text{free frequencies} & A_n/i\text{ from BG} & A_n/i\text{ from formula}\\
\hline
5 & \{2,5/2,3\} & -2304 & -2304\\
5 & \{5,1,2\} & -1760 & -1760\\
5 & \{-1,2,5\} & 14336/243 & 14336/243\\
6 & \{3/2,2,5/2,3\} & -11907/4 & -11907/4\\
6 & \{1,-2,3,4\} & -309248/2187 & -309248/2187\\
6 & \{5,1,2,3\} & -172800 & -172800\\
7 & \{3/2,2,5/2,3,7/2\} & -7302393/400 & -7302393/400\\
7 & \{1,-2,3,4,5\} & -5568/11 & -5568/11\\
7 & \{5,1,2,3,9/2\} & -9734734015/248 & -9734734015/248
\end{array}
\]
The four-point limiting checks also agreed exactly:
\[
    \{-3,2,3,-2\}\mapsto -192,
    \qquad
    \{-5,1,5,-1\}\mapsto -40.
\]

The packaged deliverables are therefore the written formula in
\texttt{answer.md}, the implementation in \texttt{two\_minus\_formula.wl}, and
the reproducibility check in \texttt{verify.wl}. The scratch files under
\texttt{codex\_work/} record the exact sample generation, symbolic chamber
calculations, failed cubic-only test, additional seven-point checks, and
four-point limiting calculation.

\section*{Expanded Public Reasoning Trace}

The rest of this document gives a more detailed reconstruction of the reasoning
process, written as an audit-friendly derivation rather than as hidden
chain-of-thought. The aim is to make clear which observations forced which
next checks, and how the formula was narrowed down.

\subsection*{1. Interpreting the computational problem}

The prompt says the answer is a piecewise homogeneous polynomial. That already
rules out several possible end states. If a direct BG evaluation produced
rational functions in a free parametrization, the denominators should cancel
after rewriting in the physical on-shell variables and after fixing a chamber.
Thus apparent rational expressions from symbolic free-frequency coordinates
were treated as coordinate artifacts, not as evidence that the final answer
should contain denominators.

The only source of chamber dependence in the supplied code is the repeated use
of
\[
    \operatorname{Abs}(k)
\]
inside the kernels and propagators. Since, in this sector,
\[
    k_1=-\omega_1^2,\qquad
    k_2=-\omega_2^2,\qquad
    k_j=\omega_j^2\quad(j\geq 3),
\]
all chamber walls should be expressible as vanishing of sums of signed squares.
This made it natural to look for inequalities involving positive-leg square
sums and the two negative square magnitudes.

\subsection*{2. Why exact rational samples came first}

The first useful question was whether the answer had a stable phase and a
stable degree. Exact rational samples showed that \(A_n/i\) was real in all
tested two-minus cases. That justified factoring out one global \(i\) and
fitting real polynomials. For example,
\[
    \{2,5/2,3\}\mapsto A_5/i=-2304,
\]
and
\[
    \{3/2,2,5/2,3\}\mapsto A_6/i=-11907/4.
\]
The scale test then showed that the five-point amplitude scales like
\(\lambda^6\). Since \(6=2\cdot5-4\), this suggested the degree \(2n-4\).
This was important because the later candidate term
\[
    \omega_1\omega_2\,(r-Q_S)^{n-3}
\]
has degree
\[
    2+2(n-3)=2n-4,
\]
exactly matching that evidence.

\subsection*{3. Why the first symbolic chamber did not determine the answer}

The first symbolic five-point chamber reduced to
\[
    \frac{A_5}{i}
    =
    \frac{-16x^5(xy+y^2+xz+yz+z^2)}{x+y+z}.
\]
Using the on-shell expression for \(\omega_1\), this can be rewritten as
\[
    \frac{A_5}{i}=16\omega_1\omega_2\,\omega_2^4.
\]
Equivalently, with
\[
    r=\min(\omega_1^2,\omega_2^2),
\]
this chamber gives
\[
    \frac{A_5}{i}=16\omega_1\omega_2\,r^2.
\]
At this point there were multiple plausible explanations:
\[
\begin{aligned}
    &G_5=r^2,\\
    &G_5\text{ is a monomial selected by chamber},\\
    &G_5\text{ is a truncated sum with only one active term}.
\end{aligned}
\]
Here
\[
    G_5=\frac{A_5}{i\,16\omega_1\omega_2}.
\]
The first chamber alone cannot distinguish these possibilities, because all
nonempty positive-leg subset sums were larger than \(r\), so every truncated
nonempty contribution would be invisible there.

\subsection*{4. The chamber that exposed finite differences}

The sample
\[
    \omega=\{-11/4,5,1,2,-21/4\}
\]
was decisive because more than one nonempty positive subset fit below the
cutoff:
\[
    r=\frac{121}{16},\qquad
    q_3=1,\qquad
    q_4=4,\qquad
    q_5=\frac{441}{16}.
\]
Terms involving \(q_5\) are inactive because \(q_5>r\). The active subsets of
\(\{3,4,5\}\) are therefore
\[
    \varnothing,\quad \{3\},\quad \{4\},\quad \{3,4\}.
\]
The normalized BG value is
\[
    G_5=\frac{-1760}{16(-55/4)}=8,
\]
using \(A_5/i=-1760\) and \(\omega_1\omega_2=-55/4\). The alternating finite
difference gives
\[
    r^2-(r-q_3)^2-(r-q_4)^2+(r-q_3-q_4)^2.
\]
Expanding this expression,
\[
\begin{aligned}
&r^2-(r-q_3)^2-(r-q_4)^2+(r-q_3-q_4)^2\\
&=r^2-(r^2-2rq_3+q_3^2)-(r^2-2rq_4+q_4^2)\\
&\qquad +(r^2-2r(q_3+q_4)+(q_3+q_4)^2)\\
&=2q_3q_4.
\end{aligned}
\]
For \(q_3=1\) and \(q_4=4\), this equals \(8\), exactly the normalized BG
value. This ruled out the naive \(G_5=r^2\) hypothesis and pointed to a
finite-difference operator.

\subsection*{5. A one-active-subset check}

The sample
\[
    \omega=\{-16/3,-1,2,5,-2/3\}
\]
has
\[
    r=1,\qquad q_3=4,\qquad q_4=25,\qquad q_5=4/9.
\]
Only the subset \(\{5\}\) is active among nonempty subsets. The finite-difference
prediction is therefore
\[
    G_5=r^2-(r-q_5)^2
    =
    1-\left(1-\frac49\right)^2
    =
    \frac{56}{81}.
\]
The BG value gives
\[
    G_5
    =
    \frac{14336/243}{16(16/3)}
    =
    \frac{14336}{243}\cdot\frac{3}{256}
    =
    \frac{56}{81}.
\]
This check tested a different chamber type: only one positive square fits below
the cutoff, rather than two whose joint subset also fits.

\subsection*{6. Why subset sums, not just individual positive legs}

The chamber conditions could not be only \(q_j<r\) for individual legs. The
finite-difference term for a subset depends on
\[
    Q_S=\sum_{j\in S}q_j.
\]
For example, if \(q_3<r\) and \(q_4<r\), the mixed term \((r-q_3-q_4)^2\) is
active only when \(q_3+q_4<r\). Thus the chamber walls must include
\[
    Q_S=r
\]
for arbitrary subsets \(S\) of positive legs. This matches the BG code's
structure: internal momenta are sums of external momenta, so a current can see a
whole subset of positive squares at once.

\subsection*{7. Rejected shortcuts}

Several simpler explanations were checked and rejected.

First, the amplitude is not just the cubic-tree part. A cubic-only recursion
gave, for the representative five-point sample,
\[
    (A_5/i)_{\mathrm{full}}=-2304,
    \qquad
    (A_5/i)_{\mathrm{cubic}}=-5296.
\]
So contact terms are essential.

Second, the five-point symbolic expressions were not best understood as
unrelated chamber monomials. Some chambers gave compact monomials after
simplification, but others gave longer expressions. The common structure only
became visible after normalizing by \(i16\omega_1\omega_2\) and rewriting in
terms of \(r\) and the \(q_j\).

Third, the four-point raw BG output could not be used naively because the real
four-point resonant locus is a boundary. The correct comparison required a
limiting calculation, not a direct call that returns \(\texttt{Indeterminate}\).

\subsection*{8. From five points to the all-\(n\) ansatz}

Once the five-point normalized result had the form
\[
    G_5(r;q_3,q_4,q_5)
    =
    \sum_{S\subseteq\{3,4,5\}}
       (-1)^{|S|}(r-Q_S)_+^2,
\]
there were two natural generalizations to test:
\[
    \sum_S(-1)^{|S|}(r-Q_S)_+^2
    \quad\text{for all }n,
\]
or
\[
    \sum_S(-1)^{|S|}(r-Q_S)_+^{n-3}.
\]
The scaling degree decides between them. Since \(r-Q_S\) is quadratic in the
frequencies, the first option would always contribute degree \(4\), and after
multiplication by \(\omega_1\omega_2\) would always have degree \(6\). That
cannot match the six- and seven-point amplitudes. The exponent must therefore
increase with \(n\), and the observed degree \(2n-4\) forces the power
\[
    n-3.
\]

The normalization also changes with \(n\). Comparing exact values at five, six,
and seven points showed the prefactor
\[
    i\,2^{n-1}\omega_1\omega_2.
\]
For \(n=5\) this is \(i16\omega_1\omega_2\), exactly the normalization used
above.

\subsection*{9. Finite-difference interpretation}

Define
\[
    H_m^{(p)}(r;q_1,\ldots,q_m)
    =
    \sum_{S\subseteq\{1,\ldots,m\}}
       (-1)^{|S|}\left(r-\sum_{j\in S}q_j\right)_+^p.
\]
Adding a new positive leg of square \(q_m\) splits the subsets into those that
do not contain \(m\) and those that do:
\[
\begin{aligned}
H_m^{(p)}(r;q_1,\ldots,q_m)
&=
H_{m-1}^{(p)}(r;q_1,\ldots,q_{m-1})\\
&\quad -
H_{m-1}^{(p)}(r-q_m;q_1,\ldots,q_{m-1}).
\end{aligned}
\]
Thus the subset formula is exactly an iterated finite difference of the
truncated monomial \(r^p\). This explains why the chamber polynomial changes by
turning terms on or off at subset-sum walls, and why alternating signs appear.

For the physical problem, \(m=n-2\) positive legs and \(p=n-3\). The final
normalized object is
\[
    G_n
    =
    H_{n-2}^{(n-3)}(r;\omega_3^2,\ldots,\omega_n^2).
\]
Multiplying by the empirically fixed universal factor gives
\[
    A_n=i\,2^{n-1}\omega_1\omega_2\,G_n.
\]

\subsection*{10. What would have falsified the formula}

The formula would have failed in several easy-to-detect ways. A chamber with
the same active subset pattern but a different polynomial value would have
ruled out the finite-difference form. A six- or seven-point sample with exact
nonzero difference from BG would have ruled out the exponent \(n-3\) or the
prefactor \(2^{n-1}\). A point where the formula had the correct absolute value
but wrong sign would have indicated a missing sign factor depending on the
frequencies. None of these failures occurred in the rational checks.

\subsection*{11. Verification strategy}

The final verification was deliberately exact rather than floating point. The
script \texttt{verify.wl} computes
\[
    \frac{A_n^{\mathrm{BG}}-A_n^{\mathrm{formula}}}{i}
\]
and simplifies the result. For all packaged non-boundary checks at
\(n=5,6,7\), this difference was exactly zero. The four-point checks used the
separate conserved-momentum limiting script and also gave exact zero difference
after taking the limit. Exact zero in rational arithmetic is stronger than the
requested \(10^{-10}\) relative numerical agreement on those test points.

\end{document}
